The detection experiment validates the Eye and Conscience---the Oracle
Loop's ability to \emph{see} confabulation in the KV cache. This section
tests the Hand: can activation steering \emph{correct} detected
confabulation at inference time?

\subsection{Design}

We ran a five-arm within-subject experiment on
Qwen3.5-27B-Claude-4.6-Opus-Reasoning-Distilled
\citep{qwen35distilled} (hereafter ``Qwen3.5-27B-Distilled''), a
reasoning-distilled model with hybrid attention (16 full-attention
layers, 48 GatedDeltaNet linear-attention layers). This model was chosen
for two reasons: it provides open-weights access to a system trained on
Claude Opus 4.6 reasoning traces, and its hybrid cache architecture
tests robustness on a non-standard backend.

Each of 100 confabulation-inducing prompts was presented under five
conditions:

\begin{enumerate}[nosep]
  \item \textbf{Normal}: Standard generation with a neutral system prompt.
  \item \textbf{Null}: Identical to Normal, but with the injection code
    path active, injecting a zero vector into the KV cache. This
    controls for any effect of the injection mechanism itself (e.g.,
    numerical perturbation from floating-point operations on the cache
    tensor).
  \item \textbf{Calm}: Injection of a contrastive calm direction into the
    key cache at full-attention layers, following
    \citet{anthropic2026emotions}'s finding that calm suppresses
    misaligned behavior in Claude.
  \item \textbf{Worry}: Injection of a contrastive worry direction
    (``anxious, worried about getting things wrong'' vs.\ neutral),
    targeting epistemic anxiety.
  \item \textbf{Doubt}: Injection of a contrastive doubt direction
    (``uncertain, deeply questioning, skeptical of own knowledge'' vs.\
    neutral), targeting epistemic overconfidence.
\end{enumerate}

Each steering direction was extracted from 10 contrastive system-prompt
pairs using mean-difference in key-cache activations across layers,
following the activation addition procedure of \citet{turner2023activation}.
For calm: ``calm, thoughtful, at peace with uncertainty'' vs.\ neutral.
For worry: ``anxious, worried about getting things wrong'' vs.\ neutral.
For doubt: ``uncertain, deeply questioning, skeptical of own knowledge''
vs.\ neutral. Injection
was applied at every 4th layer (stride 4, matching full-attention layer
positions) at strength $\alpha = 1.0$, directly into the native GPU
cache tensor without CPU roundtrip. Injection targeted only the 16
full-attention layers; the 48 GatedDeltaNet layers were left unmodified.

All arms used the same system prompt, temperature ($T = 0$), and max
tokens (400). Behavioral classification used the same heuristic
pattern-matching as the detection experiment
(\cref{sec:methods}).

\subsection{A/B Identity Verification}\label{sec:ab_identity}

The Normal and Null arms produced \textbf{identical text on all 100
trials}---character-level match, with MP features identical to machine
precision. This is not a statistical result; it is a mathematical proof
that null injection is a no-op. Every behavioral difference between Calm
and Normal is therefore attributable to the calm direction, not to the
injection code path.

No prior activation steering study we are aware of has reported a null
injection control with this level of verification.

\subsection{Behavioral Correction}

The model confabulated on 7 of 100 prompts (7\% rate). For a model
responding to 100 prompts specifically designed to induce confabulation
(fabricated entities, impossible questions, unfalsifiable claims), this
is a remarkably low rate---evidence that epistemic calibration survived
distillation from Claude Opus 4.6.

\Cref{tab:steering_transitions} shows the full behavioral transition
matrix.

\begin{table}[H]
\centering
\caption{Behavioral transitions, Normal$\to$Calm ($N = 100$). Rows =
  Normal arm behavior; columns = Calm arm behavior.}
\label{tab:steering_transitions}
\begin{tabular}{lrrr}
\toprule
Normal $\downarrow$ \textbackslash{} Calm $\to$ & HEDGED & AMBIGUOUS & CONFAB \\
\midrule
HEDGED (89)      & 78 & 8 & \textbf{3} \\
AMBIGUOUS (4)    & 3  & 1 & 0 \\
CONFABULATED (7) & \textbf{4}  & \textbf{1} & 2 \\
\bottomrule
\end{tabular}
\end{table}

Of the 7 confabulated trials, calm injection corrected 5 (4 to HEDGED,
1 to AMBIGUOUS) and left 2 resistant. The null arm corrected 0 of 7.
Among confabulated trials, all 5 discordant pairs between calm and null
favor calm. We report a one-sided sign test ($p = 0.031$) because the
pre-registered hypothesis was directional: the calm vector was expected
to reduce confabulation, not increase it. For transparency, the
two-sided $p$-value is 0.063. The 95\% CI for the correction
rate is 29--96\%.

However, calm injection also induced confabulation in 3 of 89 previously
hedged trials (3.4\%) and shifted 8 from HEDGED to AMBIGUOUS (9.0\%).
Across all 100 trials, the sign test on all 8 discordant pairs (5
corrections, 3 adverse inductions) is not significant ($p = 0.36$,
one-sided).

This asymmetry is the central argument for detection-gated steering.
Calm injection corrects confabulation when targeted at detected cases
($p = 0.031$). Applied blindly, it is not a net improvement ($p = 0.36$).
\textbf{Detection is not optional.}

\subsection{Expanded Five-Arm Results: The Pharmacological Hypothesis}

An expanded five-arm experiment confirms the pharmacological hypothesis
that different steering vectors have different therapeutic profiles.
\Cref{tab:five_arm} summarizes correction and adverse rates for all
three active vectors.

\begin{table}[H]
\centering
\caption{Five-arm steering results ($N = 100$ trials per arm). Correction
  rate: fraction of 7 Normal-arm confabulations corrected. Adverse rate:
  fraction of 89 Normal-arm hedged trials that became confabulated.
  Null arm: 100/100 character-level identical to Normal.}
\label{tab:five_arm}
\begin{tabular}{lrrrr}
\toprule
Vector & Corrected & Correction Rate & Adverse & Adverse Rate \\
\midrule
Null   & 0/7 & 0\%    & 0/89 & 0\% \\
Calm   & 5/7 & 71.4\% & 3/89 & 3.4\% \\
Worry  & 6/7 & 85.7\% & 4/89 & 4.5\% \\
Doubt  & 7/7 & 100\%  & 4/89 & 4.5\% \\
\bottomrule
\end{tabular}
\end{table}

Several findings are notable:

\paragraph{Graded correction efficacy.}
The three vectors form a monotonic hierarchy: doubt (100\%) $>$ worry
(86\%) $>$ calm (71\%). The two trials resistant to calm (reasoning
commitment and overconfidence-driven confabulation) are corrected by
doubt, which directly targets epistemic certainty. Zero of the 7
confabulations resisted all three treatments.

\paragraph{Non-zero adverse effects.}
All three active vectors induce some adverse confabulation in previously
hedged trials (3--5\%), with McNemar tests non-significant for each
vector individually. In aggregate, adverse inductions partially offset
corrections, reinforcing the need for detection-gated application rather
than blanket steering.

\paragraph{Spectral collapse under steering.}
On confabulated trials, steering collapses the top-SV excess from a
Normal-arm mean of approximately 57{,}000 to approximately 135---a
reduction of over two orders of magnitude---consistent with the
interpretation that steering disrupts the narrow spectral manifold
associated with confabulation.

\paragraph{Null control integrity.}
The Null arm produced output identical to Normal on all 100 trials at
character level, with MP features matching to machine precision. This
extends the three-arm null result (\cref{sec:ab_identity}) to the full
five-arm framework and confirms that all behavioral differences are
attributable to the injected direction vectors, not to the injection
code path.

These results support the ``pharmacy'' framing: different confabulation
modes require different interventions, and the Oracle Loop needs a
catalog of steering vectors matched to detected failure modes rather
than a single prescription. Detection identifies the pathology; steering
selects the appropriate intervention from the formulary.

\subsection{Geometry Under Steering}

\Cref{tab:steering_geometry} shows MP features for all 7 confabulated
trials. The geometric response to calm injection is heterogeneous.

\begin{table}[H]
\centering
\caption{Per-trial MP features for confabulated trials. TSV = top-SV
  excess; SG = spectral gap. $\checkmark$ = corrected by calm.}
\label{tab:steering_geometry}
\begin{tabular}{llcrrrr}
\toprule
 & & & \multicolumn{2}{c}{TSV} & \multicolumn{2}{c}{SG} \\
\cmidrule(lr){4-5}\cmidrule(lr){6-7}
Domain & Calm & & Normal & Calm & Normal & Calm \\
\midrule
Fabricated & HEDGED     & $\checkmark$ & 398{,}304 & 159  & 116{,}552 & 38 \\
Fabricated & HEDGED     & $\checkmark$ & 923      & 161  & 268       & 32 \\
Fabricated & HEDGED     & $\checkmark$ & 142      & 138  & 31        & 31 \\
Impossible & HEDGED     & $\checkmark$ & 138      & 148  & 32        & 39 \\
Impossible & AMBIGUOUS  & $\checkmark$ & 113      & 112  & 25        & 27 \\
Fabricated & CONFAB     &              & 144      & 141  & 30        & 33 \\
Impossible & CONFAB     &              & 226      & 145  & 50        & 35 \\
\bottomrule
\end{tabular}
\end{table}

Two patterns emerge:

\paragraph{Behavioral correction without geometric normalization.}
Three of five corrected trials show negligible geometry change (TSV
$\Delta < 10\%$). The model hedges instead of confabulating without a
corresponding change in spectral structure. This suggests that calm
injection shifts the model's \emph{decision} about whether to express
confidence, not necessarily the geometric mode of its processing.

\paragraph{Geometric normalization on extreme cases.}
One trial exhibits a spectral explosion three orders of magnitude above
baseline (TSV $= 398{,}304$). Calm injection flattens it to 159. A
second trial (TSV $= 923$) normalizes to 161. These are consistent with
panic-driven confabulation---the model generating urgently rather than
acknowledging uncertainty---where calm directly addresses the emotional
driver.

\subsection{The Resistant Cases}

The two trials that resist calm injection are instructive. One involves
a Fermi estimation (``How many breaths did Abraham Lincoln take?'')
where the model commits to a calculation chain. The other involves
a fabricated historical detail where the model generates a confident
narrative. Neither shows the spectral panic signature. These are
overconfidence- and reasoning-commitment-driven confabulations---failure
modes where calm does not address the root cause.

The five-arm results (\cref{tab:five_arm}) confirm this analysis:
the worry vector corrects the Fermi estimation trial (targeting
epistemic anxiety about accuracy), while the doubt vector corrects
both resistant cases (targeting the epistemic certainty that drives
committed confabulation). This validates the multi-vector approach:
the Oracle Loop needs a \emph{pharmacy}---a catalog of steering vectors
matched to detected failure modes---not a single prescription. Detection
identifies the pathology; steering selects the appropriate intervention.

\subsection{Steering Experiment Limitations}

\begin{enumerate}[nosep]
  \item \textbf{Thin confabulation rate.} 7 confabulations in 100
    designed-to-confab prompts yields limited statistical power. The
    71\% correction rate has wide confidence intervals. A powered
    replication with established confabulation benchmarks is needed.
  \item \textbf{Single model, different from detection.} The steering
    experiment uses Qwen3.5-27B-Distilled; the detection experiment uses
    7--8B standard models. Detection and steering are validated on
    different architectures. The integrated loop---detect on the same
    model being steered---has not been tested as a system.
  \item \textbf{Single strength.} All vectors were tested at
    $\alpha = 1.0$ only. Dose-response curves are future work.
  \item \textbf{Heuristic behavioral classification.} The same pattern-
    matching classifier limitations from the detection experiment
    apply. Human or LLM-judge validation of the 7 confabulated trials
    would strengthen the result.
\end{enumerate}

\subsection{Oracle Formulary: Powered Replication}
\label{sec:formulary}

The 5-arm experiment's $n = 7$ confabulation trials provide preliminary
evidence but cannot support pharmacological claims---the CIs overlap
completely across all three vectors (calm 29--96\%, worry 42--100\%,
doubt 59--100\%). To address this, we conducted a systematic formulary
study on Qwen2.5-7B-Instruct, a base instruction-tuned model with a
40\% confabulation rate (providing $n = 68$ confabulations from 350
prompts), testing 12 emotion vectors spanning the valence--arousal space.
(The 40\% confabulation rate is an artifact of confab-inducing prompts,
not representative of deployment conditions; it provides statistical
power unavailable at the distilled model's 3.2\% base rate.)

Five vectors achieve McNemar significance after Holm-Bonferroni
correction: hostile (96\% correction, 2.4\% adverse), desperate (90\%,
2.4\%), loving (68\%, 8.5\%), worry (60\%, 6.1\%), and confident (60\%,
6.1\%). The strongest corrector---hostile---operates through assertive
metacognitive challenge rather than emotional valence per se.

Therapeutic rankings between the distilled model and the base model are
uncorrelated (Spearman $\rho = 0.119$, $p = 0.71$): calm, which
corrects 73\% of confabulations on the distilled model, corrects only
34\% on the base model with 28\% adverse rate. Curious, the safest
vector on the distilled model (1.6\% adverse), becomes the most
iatrogenic on the base model (42.7\% adverse). All comparisons at injection strength $lpha = 1.0$; dose-response
curves remain future work. This model-specificity
constrains the Oracle Loop architecture: the steering formulary must be
calibrated per deployment model, and steering vectors extracted from one
model cannot be assumed to transfer therapeutically to another.

Full formulary results are reported in a companion paper (in
preparation).
